\documentclass[UTF8,11pt]{ctexart}

\usepackage[a4paper,margin=2.5cm]{geometry}
\usepackage{amsmath,amssymb}
\usepackage{booktabs}
\usepackage{graphicx}
\usepackage{xcolor}
\usepackage{hyperref}

\definecolor{stellarblue}{HTML}{2563A8}
\hypersetup{colorlinks=true,linkcolor=stellarblue,urlcolor=stellarblue}

\title{StellarLatex 中文示例}
\author{浏览器中的 \LaTeX{} 工作台}
\date{\today}

\begin{document}
\maketitle

\begin{abstract}
本文档展示常用的 \LaTeX{} 排版功能，并使用 XeTeX 在浏览器中完成编译。
源文件与生成的 PDF 都保留在本地，不会上传到服务器。
\end{abstract}

\section{文字与段落}

这是普通正文。你可以使用 \textbf{粗体}、\textit{斜体}、
\underline{下划线}，也可以用 \textcolor{stellarblue}{彩色文字}突出重点。
行内公式写作 $E=mc^2$，脚注则像这样添加\footnote{这是一个脚注示例。}。

\begin{quote}
优雅的排版，应当让读者关注内容本身，而不是排版工具。
\end{quote}

项目主页：\href{https://github.com/Arxtect/StellarLatex}{StellarLatex on GitHub}。

\section{列表}

无序列表适合并列内容：
\begin{itemize}
  \item 浏览器内运行 WebAssembly；
  \item 支持中文与数学公式；
  \item 一键生成并下载 PDF。
\end{itemize}

有序列表适合描述步骤：
\begin{enumerate}
  \item 编辑左侧的源代码；
  \item 点击“编译 PDF”；
  \item 在右侧查看结果。
\end{enumerate}

\section{数学公式}

正态分布密度函数为
\begin{equation*}
  f(x)=\frac{1}{\sigma\sqrt{2\pi}}
  \exp\left(-\frac{(x-\mu)^2}{2\sigma^2}\right).
\end{equation*}
其中 $\mu$ 是均值，$\sigma$ 是标准差。上式还展示了
分数、根号、指数和自动伸缩括号。

多行公式可以使用 \texttt{align} 环境：
\begin{align*}
  (a+b)^2 &= a^2+2ab+b^2, \\
  \sum_{k=1}^{n} k &= \frac{n(n+1)}{2}.
\end{align*}

\section{表格与图片}

下面使用 \texttt{booktabs} 绘制一个常见的三线表。
\begin{table}[htbp]
  \centering
  \caption{浏览器 LaTeX 引擎对比}
  \begin{tabular}{lcc}
    \toprule
    引擎 & 中文支持 & 输出格式 \\
    \midrule
    pdfTeX & 需要额外字体配置 & PDF \\
    XeTeX  & 原生 Unicode 支持 & PDF \\
    \bottomrule
  \end{tabular}
\end{table}

下面是一个无需外部文件的图片占位区域

实际项目中可使用
\verb|\includegraphics[width=0.8\linewidth]{image.png}| 插入图片。
\begin{figure}[htbp]
  \centering
  \fbox{\rule{0pt}{3cm}\rule{0.72\linewidth}{0pt}}
  \caption{图片占位区域}
\end{figure}

\section{代码与参考文献}

简单代码可以放入 \texttt{verbatim} 环境：
\begin{verbatim}
for item in data:
    print(item)
\end{verbatim}

文献引用通常使用 \verb|\cite| 命令，并在文末通过参考文献环境列出来源。

\begin{thebibliography}{9}
  \bibitem{knuth1984texbook}
  Donald E. Knuth.
  \textit{The TeXbook}.
  Addison-Wesley, 1984.
\end{thebibliography}

\end{document}
